\section{Introduction}\label{sec:intro}
A biologist who has sequenced a new tissue wants to know which cell types are in it. The
standard route is to cluster the cells and inspect the groups. Every clustering method used
for this has a setting that controls how finely it divides the data, and the result depends
on it strongly: on Zeisel, which has 7
annotated cell types, MICA returns between 7 and 59 clusters as that setting moves across
its own default range, and Leiden between 1 and 33. Choosing it is therefore not a
detail of tuning but the main decision, and it has to be made without knowing the answer,
because the number of cell types is usually what the experiment is trying to establish.

Published benchmarks do not make that decision. In the benchmark for scMINER
\cite{pan2025scminer}, the framework we study, every method is either given the true number
of cell types or has its setting swept and keeps only the runs that reproduce that number,
with remaining ties broken by agreement with the reference labels. The reported scores are
therefore conditional on information a user does not have. This is not specific to scMINER.
Ullmann et al. \cite{ullmann2023overoptimistic} show that tuning a clustering algorithm
against the ground truth inflates its apparent advantage, and a recent single-cell survey
\cite{szmigiel2025benchmarking} reports that on six of sixteen datasets the common default
setting performs no better than a setting picked at random. What is missing is a direct
measurement of what happens to a specific published advantage when the selection is made
without labels.

\textbf{Main question.} When the same label-free rule chooses the setting for both methods,
does MICA, the mutual information clustering in scMINER, produce partitions closer to the
known cell types than a principal component and Leiden pipeline on the same genes?

Our approach is to hold everything except the method fixed. Both methods receive the same
preprocessed matrix, choose among settings yielding 2 to 20 clusters, and are selected by
the same rule. We use a published rule, subsampling stability as specified by chooseR
\cite{pattersoncross2021chooser}, which reruns each method on 25 random 80 percent
subsamples per candidate setting and prefers settings whose clusters survive resampling. The
reference labels are used only afterwards, to score what the rule chose. The experiment
covers nine of the ten datasets in the original benchmark, repeated across three seeds, with
cost recorded.

The result is not the one we expected. Read exactly as its authors specify, the rule favours
MICA, which leads on five of eight datasets with a mean ARI of 0.675 against 0.595. Read
with a different tie-break applied to the same scores, the rule favours the baseline, 0.583
against 0.467. Nothing about the data or either method differs between those two readings.
We traced this to a measurable property: among settings the rule scores as equivalent, the
span of cluster counts averages 6.9 for MICA and 2.5 for the baseline, so for MICA the
tie-break rather than the data usually decides. On one dataset the rule cannot distinguish 3
clusters from 19. Our initial explanation, that stability would simply be less informative
for MICA, turned out to be wrong: it correlates slightly better with accuracy for MICA than
for the baseline. The signal is present but compressed into a band too narrow for a
threshold to act on.

This submission makes two contributions.

\textbf{A measurement of what label-free selection costs, and of how unstable that cost is.}
The closest work establishes that benchmarks which tune on the labels are optimistic
\cite{ullmann2023overoptimistic}, and supplies rules intended to replace that tuning
\cite{pattersoncross2021chooser,liu2021multik}, but no study applies such a rule to mutual
information clustering, and the scMINER benchmark reports no label-free result at all. We
supply that measurement and find that the uncertainty introduced by the choice of rule is
larger, on these datasets, than the difference between the two methods under any single
rule. A benchmark reporting one label-free rule could obtain either ranking.

\textbf{An account of why reproducibility fails as a criterion for this method.} MICA's
partitions are unusually reproducible under subsampling, including partitions that split a
true cell type several ways, because its distance is computed per cell pair and does not
move when the sample changes. Splitting is then as repeatable as recovering, and a rule that
rewards repeatability cannot tell them apart. We state the property that produces this,
verify it as an exact identity rather than an approximation, and use it to make the
experiment affordable, since it lets one distance matrix serve all 25 subsamples.

Section~\ref{sec:related} covers the closest work. Section~\ref{sec:method} defines the
methods and the selection rules. Section~\ref{sec:design} gives the experimental design,
Section~\ref{sec:results} the results, and Sections~\ref{sec:discussion}
and~\ref{sec:conclusion} what follows from them and what remains.
